\section{Challenges and Motivation}
\label{sec:goal}
None of the three timelines can be reconstructed from its own vantage point: the OS sees threads but not what they run, an executor sees callbacks but not where they run, and the GPU sees work but not whose it is. A complete view therefore requires all three, recorded together and tied to one another, and the architecture that decides what runs and in what order. In this section we examine what such a view requires and what stands in the way of constructing it.

\noindent
\textit{\textbf{Application Structure.}}
The application structure is the backbone of every timeline and the first thing needed to understand functional and timing behaviour: which components the software consists of, which functions are called when those components are triggered, and which channels they use to communicate. In ROS~2 the components are nodes, the functions a node manages are its callbacks, and the executor groups nodes and hands the execution of their callbacks to the OS. Knowing the structure completely then means knowing four things. Which executor serves which nodes, and how each executor is configured: single- or multi-threaded, with how many threads, and with which callback groups, which decide which callbacks may run in parallel. The kind of every dispatchable unit of work---timer, subscription, service, client or waitable---since the executor treats them differently. Every communication pattern between callbacks: producer--consumer over a topic, request--response over a service, sample-and-hold, where a callback reads the latest value another one stored, and join, where a callback fires only once a set of inputs is complete. And the mechanism behind each: a serialised message through the middleware, a pointer handed over inside the process, or state that one callback stores and another later reads, which never becomes a message and can only be inferred.

Neither source code nor existing tools deliver this structure. Source code and launch files describe what may be created, not what a given run created: objects are created and destroyed while the application runs, whether a delivery stays inside the process is decided at run time, and which callback publishes on which topic is never declared, since a publisher belongs to the node and any of its callbacks may publish through it. Tools that trace the running system recover parts. ROS-Llama~\cite{blass-2021-rosllama-paper}, Abaza et al.~\cite{abaza-2024-ros2modelsynthesis} and LiME\textsuperscript{IT}~\cite{brandenburg-2026-limeit} recover callbacks and their execution times for four of the five kinds and for single-threaded executors only; CARET and its extension~\cite{kuboichi-2022-caret,peng-2022-caretcomp} record nodes, callbacks, topics and executors on top of \texttt{ros2\_tracing}~\cite{bedard-2022-ros2tracing}, whose stock tracepoints wrap neither waitables nor the deliveries inside a process, which never reach the client library. Executor configuration is recorded by none of them.

\noindent
\textit{\textbf{CPU--GPU Interplay.}}
The CPU and the GPU are two separate domains: they are traced by different tools on different clocks, and what happens on one is invisible from the other. The GPU is moreover a scheduling domain of its own, with several engines---compute and copy---each with its own queue, and a scheduler that orders the submitted work by stream and priority, unseen by both the OS and the executor. Three things follow for a callback that offloads work. It is no longer atomic: what ROS~2 and the OS see as one execution becomes an alternating sequence of CPU segments and accelerator segments~\cite{enright-2024-paam}, and what a CPU-side trace records as its duration mixes computation with waiting for the GPU. Its segments are not fixed either: the GPU work a callback submits branches from invocation to invocation, so the sequence of segments---and hence the timing---differs between invocations of the same callback. And its GPU work must be attributed to it, where the only link between the two domains is the thread the GPU run-time was called from; under a multi-threaded executor that thread is decided at run time, so without the exact thread of every callback execution the callback--GPU relation cannot be established. No existing extractor records the accelerator at all: the ROS~2 tools above stop at the client library, and GPU profilers see kernels and copies but no callback.

\noindent
\textit{\textbf{Task Chains.}}
The executions spread over the three timelines and the two domains are not independent: together they form task chains, from a sensor reading to an actuation, and end-to-end timing can only be computed on a chain. Yet even with everything above in hand---the architecture, the communication and trigger patterns, the mapping of callbacks to nodes and executors---chains do not fall out of it. An application does not do the same thing throughout a run: some communications and executions occur only while it initialises, some only while it shuts down, and the periodic operation in between is what matters from a real-time point of view. The structure does not carry this distinction; it has to be recovered from what actually executed, and the phases separated before anything in the steady state is measured. Nor does a chain announce itself: in a deployment of hundreds of nodes the graph is entangled, a chain crosses executors, nodes and processes, and two of its links---a join, and a value sampled from stored state---leave no message and appear only as patterns in the trace. Existing tools leave the chain to the user: CARET measures a path the user names, and the extractors above produce a callback graph without the phases, the joins or the GPU segments a chain runs through.

What follows from the three challenges is the shape of the solution. The structure supplies the objects; the two domains supply the OS and GPU timelines and the segments of a callback; the chains are what runs across all of them. What is needed is one model that holds all three and the relations that tie them together---which thread ran which callback, which GPU work that callback submitted, which callback's output became which callback's input---recorded on one clock and separated by phase, so that any timeline can be reconstructed alone and all of them in combination. Such a model serves more than analysis: it lets a traced run be revisited at any granularity, and it lets the same application be re-examined under a platform the trace did not run on---another core count, scheduler, executor configuration or GPU. Section~\ref{sec:method} describes how FISH builds it.
